\chapter{Fixed Non-Commuting Damping and Modal Coupling}

The second level of the hierarchy keeps the damping matrix fixed but drops the
commutation condition \([\Lt,\Dt]=0\). The eigenvalue problem remains a QEP,
but graph modes are no longer independent. This chapter develops the resulting
modal coupling and shows that it already has the algebraic form of a rational
correction. The term non-proportional is used only in this precise sense:
damping is non-proportional for the present modal analysis when it is not
simultaneously diagonalizable with stiffness in inertia-normalized
coordinates.

\section{Modal Damping Matrix}

In graph modal coordinates,
\begin{equation}
  \ddot z+\Gamma\dot z+\Lambda z=0,\qquad
  \Gamma=Q^T\Dt Q.
  \label{eq:nonprop_modal_system}
\end{equation}
Write
\begin{equation}
  \Gamma=\Gamma_0+E,\qquad
  \Gamma_0=\diag(\Gamma_{11},\ldots,\Gamma_{nn}),
  \label{eq:nonprop_split}
\end{equation}
where \(E\) contains the off-diagonal entries. If \(E=0\), modes are decoupled.
If \(E\ne0\), the velocity of one graph mode enters the equation of another.
Equivalently, for a simple spectrum of \(\Lt\), nonzero off-diagonal entries
of \(\Gamma\) are the modal expression of the commutator \([\Lt,\Dt]\).

The off-diagonal entries are not numerical noise. They are the representation
of damping heterogeneity in the stiffness eigenbasis. A bus with large local
damping can still be ineffective for a critical mode if that mode has small
shape at the bus. Conversely, a small local damping change can matter if it
projects strongly onto the left-right structure of a sensitive eigenpair.

\section{Two-Mode Illustration}

For two graph modes, the QEP is
\begin{equation}
  \begin{bmatrix}
    s^2+s\gamma_1+\nu_1 & s\epsilon\\
    s\epsilon & s^2+s\gamma_2+\nu_2
  \end{bmatrix}
  \begin{bmatrix} z_1\\z_2\end{bmatrix}=0.
  \label{eq:two_mode_nonprop}
\end{equation}
The characteristic equation is
\begin{equation}
  (s^2+s\gamma_1+\nu_1)(s^2+s\gamma_2+\nu_2)-s^2\epsilon^2=0.
  \label{eq:two_mode_char}
\end{equation}
Eliminating \(z_2\) yields an exact scalar nonlinear equation for \(z_1\):
\begin{equation}
  s^2+s\gamma_1+\nu_1
  -
  \frac{s^2\epsilon^2}{s^2+s\gamma_2+\nu_2}
  =0.
  \label{eq:two_mode_rational}
\end{equation}
The correction is rational in \(s\). Even before converter controls are
introduced, eliminating coupled modal coordinates converts a fixed-coefficient
QEP into a frequency-dependent scalar equation for the retained mode.

\section{General Modal Correction}

Partition the modal coordinates into one target mode \(k\) and the remaining
modes \(\bar k\). The modal QEP can be written as
\begin{equation}
  \begin{bmatrix}
    p_k(s) & sE_{k\bar k}\\
    sE_{\bar k k} &
    s^2I+s\Gamma_{\bar k\bar k}+\Lambda_{\bar k\bar k}
  \end{bmatrix}
  \begin{bmatrix}
    z_k\\z_{\bar k}
  \end{bmatrix}=0,
  \label{eq:general_modal_partition}
\end{equation}
where
\begin{equation}
  p_k(s)=s^2+s\Gamma_{kk}+\nu_k.
\end{equation}
Eliminating \(z_{\bar k}\) gives
\begin{equation}
  p_k(s)
  -
  s^2E_{k\bar k}
  \left(s^2I+s\Gamma_{\bar k\bar k}+\Lambda_{\bar k\bar k}\right)^{-1}
  E_{\bar k k}
  =0.
  \label{eq:general_modal_correction}
\end{equation}
This expression shows three facts.

\begin{enumerate}
\item Off-diagonal damping affects a mode through products of coupling terms,
not through diagonal damping alone.
\item Nearby modal denominators matter more than distant denominators.
\item The correction has the same resolvent structure that later appears in
controller condensation.
\end{enumerate}

\section{Perturbation View}

For small \(E\), one may evaluate the correction near an unperturbed root
\(s_{k,0}\) of \(p_k(s)\). The leading eigenvalue shift is governed by
\begin{equation}
  \Delta s_k
  \approx
  \frac{
  s_{k,0}^2E_{k\bar k}
  \left(s_{k,0}^2I+s_{k,0}\Gamma_{\bar k\bar k}
  +\Lambda_{\bar k\bar k}\right)^{-1}
  E_{\bar k k}
  }{
  2s_{k,0}+\Gamma_{kk}
  }.
  \label{eq:nonprop_shift}
\end{equation}
The denominator \(2s_{k,0}+\Gamma_{kk}\) is the derivative of the decoupled
scalar polynomial. This is the scalar version of the general QEP sensitivity
formula. The modal correction can move both real and imaginary parts, so its
effect on damping ratio must be computed rather than inferred from graph
frequency alone.

\section{Non-Proportionality as a Gate}

The previous equations suggest a theoretical gate before using graph metrics as
damping screens. A fixed-\(D\) graph screen is plausible when
\begin{equation}
  \norm{E}\ll\norm{\Gamma_0}
  \label{eq:small_nonprop_gate_1}
\end{equation}
and when the resolvent denominators in
\eqref{eq:general_modal_correction} are not near resonance over the frequency
band of interest. This does not prove that the screen is accurate; it states
the conditions under which the modal coupling neglected by the screen is small.

Structural-dynamics non-proportionality indices formalize similar ideas:
measure how far the damping operator is from the modal subspace in which the
undamped system is diagonal \cite{conde2021nonproportionality}. The power-grid
interpretation is more specific. The off-diagonal damping matrix identifies
which graph modes exchange damping through heterogeneous device locations and
damping-to-inertia ratios.

\section{Connection to the Next Level}

Fixed non-commuting damping is not the same as converter self-energy. In
\eqref{eq:general_modal_correction}, the eliminated coordinates are other
graph modes of the same fixed-QEP system. In controller self-energy, the
eliminated coordinates are physical control and auxiliary states. The algebraic
similarity is nevertheless important. It shows that rational dependence is not
an exotic artifact; it is the natural result of eliminating coupled dynamic
coordinates. The next chapter applies that logic to controller-state
condensation.
